\section{An exact rational certificate for the mixed overlaps}
\label{sec:exact-mixed-overlap}

This appendix gives an interval-free replacement for the two mixed-overlap
signs used in Theorem~\ref{thm:witness-enclosure}.  It uses exact rational
function algebra followed by exact integer Bernstein coefficients.  The
outward-rounded certificate in the preceding appendix remains a redundant
independent check, but is not needed for this proof.

Retain the notation of Section~\ref{sec:ab-core}.  Thus
\[
 p=1-b,\qquad q=1-a,\qquad
 s=p+q,\qquad m=\min\{p,q\},\qquad M=\max\{p,q\},
\]
\[
 \chi=s^4-s^2+mM,\qquad h=\frac{\sqrt3}{2},
 \qquad c_*=c_{\max}(p,q).
\]
The strict domain gives
\begin{equation}
 mM>(1-s)(2-s),\qquad m>1-s,\qquad s>\frac23.
 \label{eq:exact-mixed-domain}
\end{equation}
Put
\[
 F_m(x)=x^4-x^2+mx-m^2,\qquad
 \kappa=F_m'(s)=4s^3-2s+m.
\]
Since $m>1-s$,
\begin{equation}
 \kappa>4s^3-3s+1=(s+1)(2s-1)^2>0.
 \label{eq:kappa-positive}
\end{equation}

\subsection{Rational upper envelopes}

Define
\begin{equation}
 \bar c_L=s-\frac{\chi}{\kappa}\quad(\chi\le0),
 \qquad
 \bar c_T=s-\frac{\chi}{1-m}\quad(\chi\ge0).
 \label{eq:branchwise-cbar}
\end{equation}

\begin{lemma}[Branchwise radial upper bounds]
\label{lem:branchwise-cbar}
On the corresponding radial cells,
\[
 c_*\le\bar c_L<1,\qquad c_*\le\bar c_T<1.
\]
At $\chi=0$ both upper bounds equal $s=c_*$.
\end{lemma}

\begin{proof}
On the $C_L(m)$ cell, $F_m(s)=\chi\le0$ and $F_m(c_*)=0$.
For $x\ge s>2/3$ one has $F_m''(x)=12x^2-2>0$, so convexity and
\eqref{eq:kappa-positive} give
\[
 0=F_m(c_*)\ge F_m(s)+F_m'(s)(c_*-s),
\]
whence $c_*\le s-\chi/\kappa=\bar c_L$.
Let $U=1-s$.  The first inequality in
\eqref{eq:exact-mixed-domain} gives $mM>U(1+U)$, and therefore
\[
 \begin{aligned}
 \chi+U\kappa
 &>U\{m+3s^3-s^2-3s+2\}\\
 &=U\{m+1-4U+8U^2-3U^3\}>0.
 \end{aligned}
\]
Indeed $0<U<1/3$ and
$1-4U+8U^2-3U^3>1-4U+7U^2>0$.  Thus $\bar c_L<1$.

On the other cell put
\[
 E=\sqrt{4s^2-3},\qquad
 m_0=\frac{s(1-E)}2,\qquad M_0=\frac{s(1+E)}2.
\]
Then
\[
 \chi=(m-m_0)(M_0-m),\qquad
 s-c_*=rac{2(m-m_0)}{1+E}.
\]
Hence
\[
 \frac{\chi}{s-c_*}
 =\frac{1+E}{2}(M_0-m)\le1-m,
\]
because $0\le E\le1$ and $M_0\le1$.  This proves
$c_*\le\bar c_T$, while $\chi>0$ gives $\bar c_T<s<1$.
\end{proof}

Reflect so that $z=a-b\ge0$ and put $U=a+b-1$.  Then
\[
 s=1-U,\qquad m=\frac{1-U-z}{2},\qquad
 D^2=z^2+6U+3U^2=:K(U,z),
\]
\[
 \Xi(U,z):=4\chi
 =1-10U+21U^2-16U^3+4U^4-z^2,
\]
\[
 G(U,z):=2\kappa=5-21U+24U^2-8U^3-z.
\]
Consequently
\begin{equation}
 \bar c_L=1-U-\frac{\Xi(U,z)}{2G(U,z)},
 \qquad
 \bar c_T=1-U-\frac{\Xi(U,z)}{2(1+U+z)}.
 \label{eq:rational-cbar}
\end{equation}
Thus both upper bounds are rational in $U,z$.

\subsection{The Gram residuals}

On the active cell write $\bar c$ for the corresponding expression in
\eqref{eq:rational-cbar}, and put
\[
 \bar\eta=h(1-\bar c),\qquad
 e=1-\bar c,\qquad t=2\bar c-1.
\]
Lemma~\ref{lem:branchwise-cbar} gives $0<\bar\eta\le\eta$ in the region
$c_*>2/3$.  Hence the three disks obtained from
$\mathbb D(C,2\eta)$, $\mathbb D(W,\eta)$, and
$\mathbb D(RA,2\eta)$ by replacing $\eta$ with $\bar\eta$ are
subsets of the disks already contained in the Minkowski sum.

Represent a global vector as $(x,hy)$ and put $C'=R^{-1}C$,
\[
 N_A=\lVert A\rVert^2,\qquad N_C=\lVert C\rVert^2,\qquad
 \nu=\langle A,C'\rangle,\qquad
 \Delta_0=x_Ay_{C'}-y_Ax_{C'}.
\]
The Gram identity is
\[
 h^2\Delta_0^2=N_AN_C-\nu^2.
\]
The two mixed square residuals for the smaller disks factor as
\begin{equation}
 \bar P_A=h^2N_AQ_A,\qquad
 Q_A=\Delta_0^2-\lVert tA-eC'\rVert^2,
 \label{eq:QA-factor}
\end{equation}
\begin{equation}
 \bar P_C=h^2N_CQ_C,\qquad
 Q_C=\Delta_0^2-\lVert tC'-eA\rVert^2.
 \label{eq:QC-factor}
\end{equation}
For instance,
\[
 \begin{aligned}
 \frac{\bar P_A}{h^2}
 &=(N_A-h^2e^2)\Delta_0^2-(tN_A-e\nu)^2\\
 &=N_A\{\Delta_0^2-t^2N_A-e^2N_C+2te\nu\},
 \end{aligned}
\]
which is \eqref{eq:QA-factor}.  Thus it remains to prove $Q_A,Q_C\ge0$.

\subsection{Exact Bernstein certificate}

The tangent witnesses are rational functions of $U,z,D$.  Substitute
\eqref{eq:rational-cbar} into \eqref{eq:QA-factor}--\eqref{eq:QC-factor}
and carry out all operations in $\mathbb Q(U,z,D)$.  Each denominator is a
positive rational constant times $(D+3)^4$, the square of the relevant
upper-envelope denominator, and the squares of the two Newton-derivative
numerators.  The latter derivatives are strictly negative by the fixed-line
sign lemma; hence the denominators are positive throughout the strict domain.

After reducing the numerator modulo $D^2-K(U,z)$, exact division gives, for
$B\in\{L,T\}$ and $X\in\{A,C\}$,
\begin{equation}
 \operatorname{num}(Q_{B,X})
 =32(K+3)^2\{A_{B,X}(U,z)+D B_{B,X}(U,z)\},
 \label{eq:exact-core-reduction}
\end{equation}
where $A_{B,X},B_{B,X}\in\mathbb Z[U,z]$.  Their term counts are
\[
\begin{array}{c|cc}
(B,X)&\#A_{B,X}&\#B_{B,X}\\ \hline
(L,A)&347&298\\
(L,C)&347&298\\
(T,A)&342&294\\
(T,C)&341&294.
\end{array}
\]

The reflected strict domain is contained in
\[
 0\le U\le\frac4{25},\qquad 0\le z\le1,
 \qquad 1-K(U,z)>0,
\]
with $\Xi\le0$ on the $L$ cell and $\Xi\ge0$ on the closure of the $T$ cell.
The verifier converts all eight polynomials in
\eqref{eq:exact-core-reduction}, together with $1-K$ and $\Xi$, to the
bivariate Bernstein basis on this rational rectangle.  It clears one common
positive denominator, and every later midpoint de Casteljau subdivision
clears a common power of two.  Therefore all stored coefficients are exact
integers.

A box is discarded only when exact Bernstein coefficients prove it lies
outside $K<1$ or outside the selected radial cell.  It is certified when every
Bernstein coefficient of the four relevant target polynomials is
nonnegative.  The convex-hull property of the Bernstein basis then proves all
four polynomials nonnegative on the box.  The deterministic subdivision has
no unresolved boxes:
\[
\begin{array}{c|rrrrrr}
&\text{processed}&\text{split}&\text{certified}&K\text{-pruned}
&\chi\text{-pruned}&\text{max depth}\\ \hline
L&165&82&43&3&37&19\\
T&191&95&43&4&49&19.
\end{array}
\]
The transcript SHA-256 values are
\[
\begin{aligned}
L:\;&\texttt{a9f738749a66b74afd19a8ba4fa7f3c8936d5e3428f040ef9bf77d121efc647d},\\
T:\;&\texttt{9775745d5147da389f1bf8d1939389cd25e600fbd7b1abcbfcdfd1d691cbd6c1}.
\end{aligned}
\]
The complete verifier is
\begin{verbatim}
proof/3XXX_CE0/31XX_Nplus1/310X_all_Vd0/
3103X_residual_core/3103X_computation/
verify_rational_mixed_overlap.py
\end{verbatim}
It uses no floating-point operation, directed rounding, or interval
arithmetic.

Since $D\ge0$, nonnegativity of the eight integer polynomials proves
$Q_A,Q_C\ge0$ for $z\ge0$ by
\eqref{eq:exact-core-reduction}; reflection supplies $z\le0$.  Equations
\eqref{eq:QA-factor}--\eqref{eq:QC-factor} give the two mixed cap overlaps for
the smaller disks, hence for the original disks.  Together with the two
analytic adjacent overlaps, this completes the cap chain and proves
Theorem~\ref{thm:witness-enclosure} without interval arithmetic.
